\chapter{Excessive Worry}
\index{Excessive Worry}

\section*{Definition and Overview}
Worry is a normal part of life, but excessive worry is persistent, difficult to control, and out of proportion to the situation. When worry becomes constant and interferes with daily life, sleep, and wellbeing, it may be generalised anxiety disorder (GAD) or another anxiety condition. Excessive worry affects the body as well as the mind, causing tension, fatigue, poor concentration, and irritability. It is very treatable, and with psychological therapy, lifestyle measures, and sometimes medicines, most people gain good control of their worry.

\section*{Epidemiology}
Anxiety disorders are among the most common mental health conditions, affecting about one in seven people in any year, and generalised anxiety disorder is a major contributor. Excessive worry is more common in women and often begins in early adulthood, although it can start at any age. Many people with excessive worry do not seek help, either because they see it as part of their personality or because they do not know that effective treatment exists.

\section*{Aetiology and Pathophysiology}
Excessive worry arises from a combination of biological and psychological factors.
\begin{itemize}
    \item \textbf{Genetics:} anxiety runs in families
    \item \textbf{Brain chemistry:} differences in the brain's stress and fear circuits contribute
    \item \textbf{Personality:} people who are naturally prone to negative thinking and perfectionism are more vulnerable
    \item \textbf{Life experiences:} stress, trauma, and difficult life events can trigger or worsen worry
    \item \textbf{Health problems:} chronic illness and some medicines can cause anxiety symptoms
    \item \textbf{Maintaining factors:} avoidance and reassurance-seeking keep worry going, because the person never learns that things will be okay
    \item \textbf{Physical effects:} chronic anxiety keeps the body in a state of tension, causing fatigue, muscle aches, and sleep problems
\end{itemize}

\section*{Clinical Features}
Excessive worry has mental and physical features.
\begin{itemize}
    \item Worrying about many different things, often for most days
    \item Difficulty controlling the worry
    \item Restlessness, feeling on edge, and irritability
    \item Fatigue and difficulty concentrating
    \item Muscle tension, headaches, and aches
    \item Sleep disturbance, including difficulty falling or staying asleep
    \item Avoiding situations that trigger worry
    \item Reassurance-seeking and checking
    \item Physical symptoms such as palpitations, sweating, and nausea
\end{itemize}

\section*{Differential Diagnosis}
Other conditions can cause or mimic excessive worry.
\begin{itemize}
    \item Depression, which often coexists with anxiety
    \item Panic disorder, with sudden panic attacks
    \item Health anxiety, focused on illness
    \item Obsessive-compulsive disorder
    \item Post-traumatic stress disorder
    \item Overactive thyroid and other medical causes of anxiety symptoms
    \item Medication side effects, including caffeine and stimulants
\end{itemize}

\section*{When to Seek Medical Care}
People should seek help when worry is persistent, difficult to control, or affecting sleep, work, or relationships. A GP can assess, treat, and refer.

\textbf{Contact your local emergency services for:}
\begin{itemize}
    \item Thoughts of self-harm or suicide
    \item Panic with chest pain and breathlessness that does not settle
    \item Severe anxiety with confusion or agitation
    \item Collapse or loss of consciousness
\end{itemize}

\section*{Diagnosis and Investigations}
Assessment distinguishes excessive worry from other conditions.
\begin{itemize}
    \item \textbf{History:} the nature, duration, and impact of the worry, and any triggers
    \item \textbf{Questionnaires:} validated scales measure anxiety severity
    \item \textbf{Physical assessment:} checking for medical causes, including thyroid disease
    \item \textbf{Mental health assessment:} screening for depression and other conditions
    \item \textbf{Assessment of risk:} checking for thoughts of self-harm
\end{itemize}

\section*{Management}
Management is effective and combines several approaches.
\begin{itemize}
    \item \textbf{Cognitive behavioural therapy:} the most effective treatment, teaching skills to challenge worry and reduce avoidance
    \item \textbf{Other therapies:} acceptance and commitment therapy and mindfulness
    \item \textbf{Relaxation techniques:} breathing, progressive muscle relaxation, and meditation
    \item \textbf{Problem-solving:} structured approaches to the real problems behind the worry
    \item \textbf{Medicines:} antidepressants and sometimes short-term anxiety medicines
    \item \textbf{Lifestyle:} regular exercise, good sleep, and reduced caffeine and alcohol
    \item \textbf{Support:} peer support and family involvement
\end{itemize}

\section*{Complications}
\begin{itemize}
    \item Depression, which commonly follows chronic anxiety
    \item Sleep problems and fatigue
    \item Physical effects of chronic stress, including high blood pressure
    \item Reduced work performance and absenteeism
    \item Relationship strain
    \item Avoidance that shrinks the person's world
\end{itemize}

\section*{Prognosis and Outcomes}
Excessive worry responds well to treatment, and most people improve substantially with cognitive behavioural therapy, alone or with medicines. The condition tends to be chronic without treatment, but the outlook with treatment is good, and many people learn to manage their worry effectively and regain their quality of life.

\section*{Prevention and Self-Care}
\begin{itemize}
    \item Set aside a limited "worry time" and postpone worrying until then
    \item Challenge worry with questions: is this worry realistic? what can I control?
    \item Practise relaxation and breathing exercises daily
    \item Exercise regularly and keep a regular sleep routine
    \item Limit caffeine and alcohol
    \item Stay connected with friends and family
    \item Seek help early from a GP
\end{itemize}

\section*{Sources and Further Reading}
The following reputable sources provide further information for healthcare students and patients seeking reliable, current advice about excessive worry and anxiety.
\begin{itemize}
    \item Beyond Blue. \textit{Anxiety information and support.} https://www.beyondblue.org.au/
    \item Head to Health. \textit{Anxiety resources.} https://www.headtohealth.gov.au/
    \item Black Dog Institute. \textit{Anxiety and worry.} https://www.blackdoginstitute.org.au/
    \item NHS. \textit{Generalised anxiety disorder.} https://www.nhs.uk/mental-health/conditions/generalised-anxiety-disorder/
    \item Mayo Clinic. \textit{Generalized anxiety disorder.} https://www.mayoclinic.org/diseases-conditions/generalized-anxiety-disorder/
    \item Australian Psychological Society. \textit{Finding a psychologist.} https://psychology.org.au/
\end{itemize}
